\clearpage\thispagestyle{empty}
\pdfbookmark[0]{Quellen zum Brief an Bode vom 3. März 1803}{bode18030303-sources}
\begin{center}\textbf{Quellen zum Brief an Bode vom 3. März 1803}\end{center}
\noindent Diese KI-Übertragung in heutiges Deutsch folgt Carl Friedrich Gauß’ Brief an Johann Elert Bode vom 3. März 1803 in \href{https://gdz.sub.uni-goettingen.de/id/PPN236007467}{Gauß’ \textit{Werken}, Band~VI} (Göttingen, 1874), S.~234--235. Gauß vergleicht Olbers’ Beobachtungen nach der Wiederauffindung der Pallas mit den Elementen für die \href{https://api.digitale-sammlungen.de/iiif/image/v2/bsb10538598_00597/full/1450,/0/default.jpg}{Ephemeride vom Dezember 1802}. Anschließend gibt er vorläufig verbesserte Elemente und einen neuen Vergleich an.
\par\smallskip\noindent Der Erstdruck im \href{https://api.digitale-sammlungen.de/iiif/image/v2/bsb10538335_00189/full/1450,/0/default.jpg}{\textit{Astronomischen Jahrbuch für 1806}, S.~179} und \href{https://api.digitale-sammlungen.de/iiif/image/v2/bsb10538335_00190/full/1450,/0/default.jpg}{S.~180} (Berlin, 1803), nennt Pallas im Titel. In seiner ersten Vergleichstafel fehlen die mittleren Bremer Zeiten. Seine Fußnote verweist auf das \href{https://api.digitale-sammlungen.de/iiif/image/v2/bsb10538334_00238/full/1450,/0/default.jpg}{\textit{Astronomische Jahrbuch für 1805}, S.~228}. Der Brief in den \textit{Werken} enthält einen Verweis in eckigen Klammern auf S.~233 des Bandes.
\par\smallskip\noindent Übertragung und Satz wurden mit KI erstellt. Die Digitalisate der Jahrbücher und der Zeitschrift stammen von der Bayerischen Staatsbibliothek.
\clearpage
